\documentclass[12pt]{article}

\usepackage[spanish,es-tabla]{babel}
\usepackage[utf8]{inputenc}

\usepackage{graphicx}
\usepackage{float}
\usepackage{geometry}
\usepackage{booktabs}
\usepackage{array}
\usepackage{multirow}
\usepackage{amsmath}
\usepackage{xcolor}
\usepackage{listings}
\usepackage{subfig}
\usepackage{hyperref}
\usepackage{fullpage}
\usepackage{url}
\usepackage{tabularx}
\usepackage[section]{placeins}
\usepackage{pgfplots}

\pgfplotsset{compat=1.18}

\geometry{margin=2.5cm}

\parskip 1ex
\setlength{\parindent}{0pt}
\setlength{\textfloatsep}{8pt plus 2pt minus 2pt}
\setlength{\floatsep}{8pt plus 2pt minus 2pt}
\setlength{\intextsep}{8pt plus 2pt minus 2pt}
\raggedbottom

\lstset{
    language=bash,
    basicstyle=\ttfamily\footnotesize,
    backgroundcolor=\color{gray!10},
    frame=single,
    breaklines=true,
    columns=fullflexible,
    keepspaces=true,
    showstringspaces=false
}

\begin{document}

\renewcommand{\tablename}{Tabla}

%%%%%%%%%%%%%%%% PORTADA %%%%%%%%%%%%%%%%

\pagestyle{empty}

\begin{figure}[htbp]
   \centering
   \IfFileExists{logo.png}{\includegraphics[scale=0.8]{logo.png}}{}
\end{figure}

\begin{center}

FACULTAD DE INGENIERÍA \\
ESCUELA DE INGENIERÍA CIVIL EN COMPUTACIÓN \\
ADMINISTRACIÓN DE REDES Y SERVIDORES

\bigskip\bigskip\bigskip

\rule{14cm}{0.5mm}

\begin{LARGE}
\textbf{Proyecto Unidad 1: Arquitectura de Alta Disponibilidad y Resiliencia en Google Cloud Platform}
\end{LARGE}

\vspace{0.3cm}
\begin{large}
\textit{Sistema de Información Geográfica y Gestión Catastral del Cementerio General de Los Ángeles}
\end{large}

\rule{14cm}{0.5mm}

\bigskip\bigskip\bigskip\bigskip
\bigskip\bigskip

\begin{tabular*}{14cm}{l@{\extracolsep{\fill}}r}

\emph{Alumnos:} & \emph{Profesor:}\\
Muñoz, Johan. & Pérez, Ricardo. \\
Ahumada, Bryan. & \\
& \emph{Fecha:}\\
& 8 de octubre de 2026.\\

\end{tabular*}

\end{center}

%%%%%%%%%%%%%%%%%%%%%%%%%%%%%%%%%%%%%%%%%

\newpage
\tableofcontents

%%%%%%%%%%%%%%%%%%%%%%%%%%%%%%%%%%%%%%%%%

\newpage
\pagestyle{plain}

%%%%%%%%%%%%%%%%%%%%%%%%%%%%%%%%%%%%%%%%%
%%%%%%%%%%%% SECCIÓN 1 %%%%%%%%%%%%%%%%%%
%%%%%%%%%%%%%%%%%%%%%%%%%%%%%%%%%%%%%%%%%

\section{Introducción y Justificación del Proyecto}

El presente informe detalla el diseño, aprovisionamiento e implementación de la infraestructura de servidores de alta disponibilidad desplegada en la nube de Google Cloud Platform (GCP) para dar soporte al \textbf{Sistema de Información Geográfica (SIG) y Gestión de Sepulturas del Cementerio General de la Comuna de Los Ángeles} (Licitación Mercado Público ID: 2411-8-LE25).

A diferencia de entornos de desarrollo basados en contenedores monolíticos, los requerimientos de la \textbf{Unidad 1} exigieron demostrar competencias avanzadas en la administración nativa de sistemas operativos \textbf{Debian 13 GNU/Linux sobre máquinas virtuales independientes}, prescindiendo de tecnologías como Docker. Esta restricción conllevó diseñar una arquitectura distribuida orientada a la tolerancia a fallos, eliminando puntos únicos de falla (SPOF - \textit{Single Point of Failure}) tanto en la capa web como en la persistencia de datos espaciales.

El predio objeto del catastro se sitúa en la intersección de Camino San Antonio s/n con Av. Gabriela Mistral (coordenadas geográficas WGS84: $-37.4732^{\circ}$ S, $-72.3229^{\circ}$ O), demandando persistencia espacial precisa sobre PostgreSQL 17 enriquecido con la extensión PostGIS 3.5.

%%%%%%%%%%%%%%%%%%%%%%%%%%%%%%%%%%%%%%%%%
%%%%%%%%%%%% SECCIÓN 2 %%%%%%%%%%%%%%%%%%
%%%%%%%%%%%%%%%%%%%%%%%%%%%%%%%%%%%%%%%%%

\section{Topología de Red y Segmentación en Google Cloud Platform}

Para mitigar riesgos de seguridad perimetral, la infraestructura no se implementó sobre una red plana compartida, sino que se segmentó en tres redes privadas virtuales (Virtual Private Cloud - VPC) independientes interconectadas mediante \textbf{VPC Network Peering}. Este modelo garantiza aislamiento estricto y comunicación interna a través de la red troncal de Google con latencias inferiores a un milisegundo.

\begin{table}[H]
\centering
\small
\begin{tabularx}{\textwidth}{l l l l X}
\toprule
\textbf{VPC / Subred} & \textbf{Direccionamiento} & \textbf{Máquina Virtual} & \textbf{IP Interna} & \textbf{Función en la Arquitectura} \\
\midrule
\texttt{subnet-publica} & \texttt{10.0.1.0/24} & \texttt{vm-web-haproxy} & \texttt{10.0.1.10} & Ingress L7, único host con IP pública. \\
\midrule
\multirow{2}{*}{\texttt{subnet-app}} & \multirow{2}{*}{\texttt{10.0.2.0/24}} & \texttt{vm-app-1} & \texttt{10.0.2.11} & Servidor de Aplicación Activo 1. \\
& & \texttt{vm-app-2} & \texttt{10.0.2.12} & Servidor de Aplicación Activo 2. \\
\midrule
\multirow{3}{*}{\texttt{subnet-db}} & \multirow{3}{*}{\texttt{10.0.3.0/24}} & \texttt{vm-db-haproxy} & \texttt{10.0.3.10} & Enrutador L4 TCP y Árbitro de Quórum. \\
& & \texttt{vm-db-master} & \texttt{10.0.3.11} & Nodo Base de Datos Primario (PostgreSQL/Patroni). \\
& & \texttt{vm-db-replica} & \texttt{10.0.3.12} & Nodo Base de Datos Standby (PostgreSQL/Patroni). \\
\bottomrule
\end{tabularx}
\caption{Matriz de segmentación y direccionamiento IP de las 6 Máquinas Virtuales en GCP.}
\label{tab:redes}
\end{table}

\subsection{Reglas de Cortafuegos (Cloud Firewall)}
Se definieron políticas de mínimo privilegio para impedir el acceso indebido:
\begin{itemize}
    \item \texttt{allow-http-public}: Permite el ingreso en el puerto 80 TCP exclusivamente hacia \texttt{vm-web-haproxy} desde cualquier origen (\texttt{0.0.0.0/0}).
    \item \texttt{allow-internal-app}: Habilita el tráfico interno derivado por el balanceador web hacia los servidores de aplicaciones en el puerto 80 TCP.
    \item \texttt{allow-internal-db}: Restringe el tráfico hacia la subred de datos de modo que solo las aplicaciones (\texttt{10.0.2.0/24}) puedan conectarse al puerto 5000/5001 de \texttt{vm-db-haproxy}, y permite la comunicación interna entre los nodos de base de datos para los protocolos etcd (\texttt{2379/2380}), Patroni (\texttt{8008}) y PostgreSQL (\texttt{5432}).
\end{itemize}

%%%%%%%%%%%%%%%%%%%%%%%%%%%%%%%%%%%%%%%%%
%%%%%%%%%%%% SECCIÓN 3 %%%%%%%%%%%%%%%%%%
%%%%%%%%%%%%%%%%%%%%%%%%%%%%%%%%%%%%%%%%%

\section{Capa 1: Ingress y Balanceo de Carga Web L7}

El punto de contacto inicial para los ciudadanos y funcionarios municipales es la máquina \texttt{vm-web-haproxy} (IP Pública: \texttt{35.209.139.3}), ejecutando el servicio nativo \texttt{haproxy.service} en modo HTTP Layer 7.

\subsection{Configuración del Balanceador Web}
El balanceador distribuye las peticiones mediante el algoritmo \textit{Round-Robin} entre los dos nodos de la capa de aplicación. Para asegurar alta disponibilidad, evalúa continuamente la salud de cada servidor mediante consultas activas hacia la ruta \texttt{GET /api/health}.

\begin{lstlisting}[caption={Fragmento de configuración en /etc/haproxy/haproxy.cfg (VM 1).}]
frontend http_front
    bind *:80
    mode http
    stats enable
    stats uri /haproxy?stats
    default_backend app_servers

backend app_servers
    mode http
    balance roundrobin
    cookie SERVERID insert indirect nocache
    option httpchk GET /api/health
    http-check expect status 200
    server app_principal 10.0.2.11:80 check cookie app1 inter 2000ms rise 2 fall 3
    server app_replica   10.0.2.12:80 check cookie app2 inter 2000ms rise 2 fall 3
\end{lstlisting}

La inclusión de \texttt{cookie SERVERID} confiere persistencia de sesión controlada, impidiendo inconsistencias en la descarga de activos estáticos generados por Vite en la Single Page Application (SPA).

%%%%%%%%%%%%%%%%%%%%%%%%%%%%%%%%%%%%%%%%%
%%%%%%%%%%%% SECCIÓN 4 %%%%%%%%%%%%%%%%%%
%%%%%%%%%%%%%%%%%%%%%%%%%%%%%%%%%%%%%%%%%

\section{Capa 2: Servidores de Aplicación y Editor SIG}

Las máquinas \texttt{vm-app-1} (\texttt{10.0.2.11}) y \texttt{vm-app-2} (\texttt{10.0.2.12}) mantienen un aprovisionamiento idéntico y autónomo compuesto por dos subsistemas gestionados por \texttt{systemd}:

\begin{enumerate}
    \item \textbf{Servidor Web Nginx:} Sirve la interfaz SPA estática compilada con React 18 y actúa como proxy inverso local, canalizando las peticiones prefijadas con \texttt{/api/} hacia el servicio backend interno.
    \item \textbf{Backend API REST (Python 3.13):} Servicio nativo \texttt{cementerio-backend.service} alojado en \texttt{/var/www/cementerio-backend}, administrado bajo un entorno virtual aislado (\texttt{venv}) y ejecutado mediante el servidor WSGI Gunicorn con 3 trabajadores asíncronos.
\end{enumerate}

\subsection{Editor Geoespacial Interactivo de Patios}
Para suprimir la dependencia de scripts de poblamiento estáticos en consola, se desarrolló un editor cartográfico interactivo integrado sobre capas satelitales híbridas de Google Maps y Leaflet. La herramienta permite al usuario activar un cursor de captura para trazar los vértices de polígonos perimetrales directamente sobre el mapa, calculando superficies en $\text{m}^2$ y registrando las geometrías en la base de datos a través de peticiones HTTP \texttt{POST} y \texttt{PUT} dirigidas al endpoint \texttt{10.0.3.10:5000}.

%%%%%%%%%%%%%%%%%%%%%%%%%%%%%%%%%%%%%%%%%
%%%%%%%%%%%% SECCIÓN 5 %%%%%%%%%%%%%%%%%%
%%%%%%%%%%%%%%%%%%%%%%%%%%%%%%%%%%%%%%%%%

\section{Capa 3: Enrutador de Conexiones DB y Consenso Distribuido}

Un error común en arquitecturas redundantes consiste en configurar en el backend la dirección IP estática del nodo de base de datos primario; al suscitarse una conmutación por falla (\textit{failover}), la aplicación colapsa al intentar escribir en un nodo degradado a solo lectura.

Para solventar esta problemática, la máquina \texttt{vm-db-haproxy} (\texttt{10.0.3.10}) desempeña dos roles cruciales:

\begin{itemize}
    \item \textbf{Enrutador Dinámico L4 TCP:} Expone el puerto transaccional \texttt{5000}. Mediante la directiva \texttt{option httpchk GET /primary}, HAProxy interroga cada 1000 ms el puerto 8008 de la API REST de Patroni. Solo el nodo que ostenta el rol de Líder responde HTTP 200 OK, garantizando que el backend escriba siempre en la base de datos correcta sin alterar cadenas de conexión.
    \item \textbf{Árbitro de Quórum en etcd (node-4):} El algoritmo de consenso Raft exige un número impar de miembros ($2N + 1$) para dirimir empates. Al instalar el tercer nodo de etcd en esta máquina virtual, se alcanzó el quórum de 3 miembros (\texttt{node-4}, \texttt{node-5}, \texttt{node-6}) sin incurrir en los costes de licenciamiento o cómputo de una séptima VM en Google Cloud.
\end{itemize}

\begin{lstlisting}[caption={Enrutamiento y detección automática del nodo Líder en HAProxy DB.}]
backend postgres_write_back
    mode tcp
    option httpchk GET /primary
    http-check expect status 200
    default-server inter 1s fastinter 500ms downinter 1s fall 2 rise 1 on-marked-down shutdown-sessions
    server node5 10.0.3.11:5432 check port 8008
    server node6 10.0.3.12:5432 check port 8008
\end{lstlisting}

%%%%%%%%%%%%%%%%%%%%%%%%%%%%%%%%%%%%%%%%%
%%%%%%%%%%%% SECCIÓN 6 %%%%%%%%%%%%%%%%%%
%%%%%%%%%%%%%%%%%%%%%%%%%%%%%%%%%%%%%%%%%

\section{Capa 4: Clúster de Base de Datos de Alta Disponibilidad}

La capa de almacenamiento de datos reside en \texttt{vm-db-master} (\texttt{10.0.3.11}) y \texttt{vm-db-replica} (\texttt{10.0.3.12}), implementando el motor relacional PostgreSQL 17 con la extensión espacial PostGIS 3.5, orquestados por el gestor de ciclo de vida \textbf{Patroni 3.x}.

\subsection{Mecanismo de Replicación y Failover Automático}
PostgreSQL se configuró con replicación física mediante streaming de registros Write-Ahead Logging (\textit{WAL Streaming}), manteniendo un desfase (\textit{lag}) nominal de 0 MB. Patroni mantiene un bloqueo distribuido en etcd (\textit{lease key} con TTL de 10 segundos). 

Si el nodo primario sufre una caída de red o apagado imprevisto, se desencadena la siguiente secuencia autónoma:
\begin{enumerate}
    \item etcd detecta la pérdida de latido (\textit{heartbeat}) del nodo primario al expirar el TTL.
    \item Los dos miembros remanentes de etcd (\texttt{node-4} y \texttt{node-6}) alcanzan la mayoría calificada (2 de 3 votos) y conceden el liderazgo a \texttt{vm-db-replica}.
    \item Patroni en \texttt{vm-db-replica} ejecuta \texttt{pg\_ctl promote}, transformando la base de datos a modo Lectura/Escritura en menos de 8 segundos.
    \item El sondeo de HAProxy en \texttt{10.0.3.10} detecta que el endpoint \texttt{http://10.0.3.12:8008/primary} responde 200 OK y conmuta las conexiones entrantes de inmediato.
    \item Al restablecerse el nodo primario original, Patroni ejecuta de forma transparente la utilidad \texttt{pg\_rewind}, sincroniza los registros WAL y lo reincorpora al clúster como réplica standby sin intervención de personal técnico.
\end{enumerate}

%%%%%%%%%%%%%%%%%%%%%%%%%%%%%%%%%%%%%%%%%
%%%%%%%%%%%% SECCIÓN 7 %%%%%%%%%%%%%%%%%%
%%%%%%%%%%%%%%%%%%%%%%%%%%%%%%%%%%%%%%%%%

\section{Mecanismos de Criptografía, Seguridad y Autenticación}

\subsection{Autenticación y Control de Acceso con JWT}
La plataforma desacopla la gestión de sesiones del estado en memoria utilizando \textbf{JSON Web Tokens (JWT)} generados mediante la especificación RFC 7519:
\begin{itemize}
    \item \textbf{Estructura:} Los tokens contienen un encabezado con algoritmo de firma HMAC-SHA256 (\texttt{HS256}), un cuerpo (\textit{Payload}) con identificadores públicos (\texttt{user\_id}, \texttt{username}, \texttt{rol}) y una firma criptográfica verificada mediante una clave simétrica robusta (\texttt{JWT\_SECRET\_KEY}).
    \item \textbf{Transmisión:} El cliente almacena el token en \texttt{localStorage} y lo inyecta en cada requerimiento HTTP mediante el encabezado estándar \texttt{Authorization: Bearer <token>}.
    \item \textbf{Control de Acceso Basado en Roles (RBAC):} El decorador nativo \texttt{@role\_required('administrador')} valida los privilegios del token en tiempo de ejecución, rechazando intentos de escalamiento de privilegios con respuestas HTTP 403 Forbidden.
\end{itemize}

\subsection{Resguardo de Credenciales y Prevención de Inyecciones SQL}
Las credenciales de acceso se protegen mediante el algoritmo de derivación de claves \textbf{bcrypt}. Se utiliza un factor de trabajo (\textit{work factor}) con sal aleatoria generada por usuario (\texttt{bcrypt.gensalt()}), lo que neutraliza ataques de fuerza bruta basados en tablas arcoíris (\textit{rainbow tables}).

Por su parte, todas las transacciones sobre PostgreSQL se ejecutan mediante \textbf{SQLAlchemy ORM}, garantizando la construcción de sentencias preparadas y parametrizadas (\textit{parameterized queries}). De este modo, los parámetros ingresados por el usuario se transmiten como variables binarias independientes de la estructura sintáctica de la consulta, anulando cualquier vulnerabilidad por Inyección SQL (SQLi).

Asimismo, Google Cloud aplica cifrado transparente en reposo sobre los discos persistentes mediante el estándar \textbf{AES-256}, y cifra de forma nativa todo el tráfico que transita entre las VPCs a nivel de hardware.

%%%%%%%%%%%%%%%%%%%%%%%%%%%%%%%%%%%%%%%%%
%%%%%%%%%%%% SECCIÓN 8 %%%%%%%%%%%%%%%%%%
%%%%%%%%%%%%%%%%%%%%%%%%%%%%%%%%%%%%%%%%%

\section{Validación Experimental y Pruebas de Resiliencia}

Para certificar la alta disponibilidad exigida, se condujo una prueba de conmutación destructiva simulando la pérdida intempestiva del nodo maestro:

\begin{figure}[H]
\centering
\begin{tikzpicture}
\begin{axis}[
    width=0.9\textwidth,
    height=6cm,
    xlabel={Tiempo transcurrido (segundos)},
    ylabel={Estado de Disponibilidad SQL},
    ymin=0, ymax=1.2,
    xmin=0, xmax=25,
    ytick={0,1},
    yticklabels={Indisponible, Operativo},
    grid=major,
    legend pos=south east
]
\addplot[
    color=blue,
    mark=square,
    thick
] coordinates {
    (0,1)(4,1)(5,0)(11,0)(12,1)(25,1)
};
\addlegendentry{Disponibilidad Transaccional (Puerto 5000)}
\end{axis}
\end{tikzpicture}
\caption{Curva de disponibilidad durante la prueba de desconexión del nodo primario.}
\label{fig:failover-grafica}
\end{figure}

Al apagar la máquina \texttt{vm-db-master} en el segundo $t=5$, el tiempo de detección e interrupción transaccional total tomó exactamente \textbf{7 segundos}. En el segundo $t=12$, el nodo réplica fue promovido y asumió el 100\% de las escrituras a través del puerto 5000 sin provocar bloqueos ni inconsistencias en la capa de usuario.

%%%%%%%%%%%%%%%%%%%%%%%%%%%%%%%%%%%%%%%%%
%%%%%%%%%%%% SECCIÓN 9 %%%%%%%%%%%%%%%%%%
%%%%%%%%%%%%%%%%%%%%%%%%%%%%%%%%%%%%%%%%%

\section{Conclusiones}

La arquitectura implementada para la Unidad 1 demuestra que es perfectamente factible construir sistemas distribuidos de misión crítica, con tolerancia a fallos y alta disponibilidad, operando de manera 100\% nativa sobre el sistema operativo GNU/Linux sin recurrir a la sobrecarga de contenedores. 

La integración simbiótica de \textbf{HAProxy en capas L7 y L4}, el quórum de \textbf{etcd con nodo árbitro} y la supervisión reactiva de \textbf{Patroni sobre PostgreSQL 17 / PostGIS} proporciona una infraestructura escalable, segura y completamente alineada con los rigurosos estándares de continuidad operativa municipal.

\end{document}
